\documentclass[a4paper,fleqn,11pt]{article}
\usepackage{tabularx}
\usepackage{pdflscape}
\usepackage{rotating}

%%%%%%%%%%%%%%%%%%%%


\usepackage[]{geometry}
\usepackage[latin1]{inputenc}
\usepackage[UKenglish]{babel}
\usepackage[UKenglish]{isodate}
\usepackage{amsmath}
\usepackage{amsfonts}
\usepackage{amssymb}
\usepackage{amsthm}
\usepackage{graphicx}
\usepackage{chngpage}
\usepackage{calc}
\PassOptionsToPackage{hyphens}{url}
\usepackage{hyperref}
\usepackage[nameinlink]{cleveref}
\usepackage{fancyhdr}
\usepackage{titletoc}
\usepackage[explicit]{titlesec}
\usepackage{natbib}
\usepackage[dvipsnames]{xcolor}
\usepackage[sc]{mathpazo}
\linespread{1.05}
\usepackage[T1]{fontenc}
\usepackage{minted}

\hypersetup{
	colorlinks=true,
	linkcolor=black,
	urlcolor=black,
	citecolor=black
}

\setlength{\parindent}{0mm}
\setlength{\parskip}{\medskipamount}
\renewcommand\baselinestretch{1.2}

\cleanlookdateon

\makeatletter
\newcommand{\@assignment}[0]{Assignment}
\newcommand{\assignment}[1]{\renewcommand{\@assignment}{#1}}
\newcommand{\@supervisor}[0]{}
\newcommand{\supervisor}[1]{\renewcommand{\@supervisor}{#1}}
\newcommand{\@yearofstudy}[0]{}
\newcommand{\yearofstudy}[1]{\renewcommand{\@yearofstudy}{#1}}
\makeatletter

\newtoggle{IsDissertation}
%%%%%%%%%%%%%%%%%%%%%%%%%%%%%%%%%%%%%%%%%%%%%%%%%%%%%%%%%%%%%%%%%%%%%%%%%%%%%%%
%% Project-specific configuration
%%%%%%%%%%%%%%%%%%%%%%%%%%%%%%%%%%%%%%%%%%%%%%%%%%%%%%%%%%%%%%%%%%%%%%%%%%%%%%%

\author{Matthew Knight}
\title{Design and Analysis of a\\ Heterogeneous RISC-V SoC\\ on FPGAs}
\supervisor{Dr Eduardo W\"{a}chter}
% \yearofstudy{3\textsuperscript{rd}}

%%%%%%%%%%%%%%%%%%%%%%%%%%%%%%%%%%%%%%%%%%%%%%%%%%%%%%%%%%%%%%%%%%%%%%%%%%%%%%%


\assignment{Project specification}

%%%%%%%%%%%%%%%%%%%%

\pagestyle{plain}
\renewcommand{\headrulewidth}{0.0pt}

\makeatletter
\fancypagestyle{plain}{
	\fancyhf{}
	\fancyhead[R]{\textit{\@title} - \textit{\@assignment}}
    \fancyhead[L]{\textit{\@author}}
    \fancyfoot[C]{\thepage}
}
\makeatother

\newcolumntype{b}{X}
\newcolumntype{s}{>{\hsize=.5\hsize}X}

\setcitestyle{square}

%%%%%%%%%%%%%%%%%%%%

\begin{document}

\makeatletter
\begin{titlepage}

    % do not show the assignment name for the dissertation
    \iftoggle{IsDissertation}
    {
        \textbf{\Huge \@title} \\[1.5cm]
    }
    {
        \textbf{\Huge \@title} \\
        \Large \@assignment \\[1.5cm]
    }
    \Large \textbf{\@author} \\
    Department of Computer Science \\
    University of Warwick \\

    % include the supervisor and year of study if the are specified and its the disseration
    \iftoggle{IsDissertation}{
        \ifdefempty{\@supervisor}{}{
            Supervised by \@supervisor \\
        }\ifdefempty{\@yearofstudy}{}{
            Year of Study: \@yearofstudy \\
        }
    }{}

    \vfill

    % include the current date for the dissertation
    \iftoggle{IsDissertation}{\today}{}

    \begin{adjustwidth}{-\oddsidemargin-1in}{-\rightmargin}
        \centering
        \includegraphics[width=\paperwidth]{../common/line.png}
    \end{adjustwidth}

    \vspace*{-3.5cm}

\end{titlepage}
\makeatother


\pagestyle{plain}

\section{Introduction}
The aim of this project is to design a heterogeneous CPU using the RISC-V open source instruction set. Once designed, the CPU should be implementable on an FPGA and be able to run a Debian Linux distro.

Designing a CPU is a complex challenge. They need to be power efficient to reduce the cost of running the system, as well as providing adequate processing power when required. Often one is sacrificed to improve the other: for instance, energy efficiency is often reduced due in larger CPU designs. Due to the increased size and complexity, more power has to be supplied to achieve the same performance as a smaller processor.

A heterogeneous CPU attempts to solve this issue by combining dissimilar core designs; often a powerful core, or p-core, and an efficient core, or e-core. When only light processing is required, the p-core can be effectively shutdown and the e-core will do all processing, resulting in less power used. For heavy processing, the p-core is then used to increase peak performance. Depending on the exact implementation, the p-core can be used either individually or in tandem with the e-core, but both result in greater performance than just the e-core.

\subsection{Definitions}
\begin{description}
    \item[Heterogeneous CPU] A CPU that contains 2 different types of core, mismatched in performance and power consumption.
    \item[RISC-V ISA] An open-source CPU instruction set architecture.
    \item[SoC] System-on-Chip, a complete computing system within a single circuit board or design. 
    \item[FPGA] A field programmable gate array, an integrated circuit that can be reprogrammed using a hardware description language (HDL).
\end{description}

\section{Background}
\subsection{Example heterogeneous CPU: big.LITTLE}
ARM released big.LITTLE in 2011\cite{biglittle}, a mobile heterogeneous architecture. The architecture provides the low power usage needed in mobile devices the majority of the time when idle, only running tasks like checking for new messages. It also provides the high performance that can be demanded from phones when used to browse the internet, play mobile games, etc. Most big.LITTLE designs use HMP (heterogeneous multiprocessing) where all cores are available to have processes assigned to them at all times. The alternatives to this are: clustered switching, where either the big cores or the LITTLE cores are in use, and in-kernel switching, where big and LITTLE cores are paired to form a virtual core and only one performs the tasks assigned to the virtual core at any one time.

\subsection{The RISC-V ISA}
The RISC-V ISA\cite{riscvISA} is an open-source instruction set for the design of a CPU, defining the behaviour of the CPU. By defining just behaviour, the actual implementation of the instructions are left to the designer and allow for characteristics like performance, size and more to vary from CPU to CPU. In addition to this, the ISA has extensions and different base options so that a CPU can be tailored to fit an exact use case where only necessary instructions are implemented.

Some of the RISC-V implementations using correct base and extension options are also able to run versions of Linux. The RV64GC configuration is an example of this, with the Debian distribution\cite{debianriscv} offering versions that directly support systems using RISC-V processors that meet that minimum configuration. In addition to this, around 95\% of the packages available on the Debian OS can be built on a RISC-V machine, allowing such a system to be used for general computing to some extent.

The majority of use cases for RISC-V are currently embedded and highly specialised systems such as compute accelerators for tasks like artificial intelligence, neural networks and image processing. General purpose use of RISC-V is not currently widespread, and design of general purpose RISC-V processors is a developing area.

\section{Existing Solutions}
\subsection{Rocket Chip generator}
Rocket Chip\cite{rocketchip} is a generator for SoCs, capable of producing many different configurations and designs of RISC-V based SoCs for use in simulation and implementation in FPGAs and VLSI (Very-large-scale integration, creation of an IC of a design). The Rocket Chip generator creates RTL code in Chisel, a hardware description language based on Scala. To create a new SoC, a configuration is written in Scala to describe the features of the SoC, such as cores, cache, tiling (organisation) and peripheral connections. 

Cores and their organisation are the main interest of this project, and Rocket Chip provides a large range of options. Rocket Chip can generate two types of RV64GC core: BOOM\cite{zhaosonicboom} and Rocket. The BOOM cores contain out-of-order execution pipelines, allowing for more efficient processing of instructions. The Rocket cores are simpler, with in-order pipelines. Both BOOM and Rocket cores have optional FPUs, configurable branch predictors and options for extra ISA extensions. The base Rocket Chip generator allows for combinations of BOOM and Rocket cores to be implemented on a single SoC, including heterogeneous designs such as a BOOM and Rocket core together. 

\subsubsection{Chipyard}
Chipyard\cite{chipyard} is an SoC design framework, and an alternative to Rocket Chip. Chipyard provides more options for adding accelerators, as well as more options for types of host core, including non-RISC options. The framework directly allows different types of core to be implemented in one SoC and are combined together using the Constellation interconnect for highly irregular architectures.

\subsection{Other Designs}
Most existing heterogeneous solutions appear to comprise of a 'host' processor and an accelerator, where the host processor is constantly running and offloads tasks to the accelerator when needed. This is different to the aim of this project, where the cores are of equal standing and both could be used individually. A brief summary of these works can be seen below.

\begin{sidewaystable}
\begin{center}
        \begin{tabularx}{\textwidth}{ssssb}
            \hline
            Title & Core     Organisation & Application & Implementation & Comments \\
            \hline
            A RISC-V Heterogeneous SoC for Embedded Devices\cite{valenterisc}. & Heterogeneous CVA6 core\cite{zaruba2019cost} with 32-bit RISC-V core. & Linux-capable for ML/DSP workloads. & Implemented on a Xilinx VCU118 with plans for produced IC on 22nm. & Early in development stage, but has a successful heterogeneous design and is Linux-capable. Still a host-accelerator combination that is less flexible than equal cores. \\
            \hline
            Muntjac multicore RV64 processor\cite{UCAM-CL-TR-972}. & Homogeneous RV64GC multicore. & Linux-capable base to enable specialised designs. & Implemented on Xilinx Kintex 7. & Aims to be easily understood and extendable - a great starting point for development of specialised designs. Limited to only one type of core in base repository, but contains interconnect for other designs. \\
            \hline
            BlackParrot\cite{blackparrot} & Homogeneous RV64G multicore. & Linux-capable general purpose processor. & Implemented on an IC at 12nm. & A general purpose RISC-V multicore. Implicitly allows heterogeneous CPU by hiding the multicore components beneath an ISA layer and allows easy extensions by keeping the design tiny, modular and friendly. \\
            \hline
        \end{tabularx}
\end{center}
\end{sidewaystable}

\clearpage
\section{Objectives}
The overall objective is design a heterogeneous RISC-V CPU to be used in an SoC that can run Linux and use both cores to improve peak performance and power efficiency. While this is the overarching goal, it can be broken down into smaller objectives. These objectives have been labelled with each labelled as Must, Should or Could to indicate their importance and expected completion.
\begin{enumerate}
    \item Design a heterogeneous RISC-V SoC containing 2 dissimilar cores, each capable of executing at least the RV64G instruction set (Must).
    \item Execute instructions on both cores when implemented in an SoC on an FPGA (Must).
    \item Have comparable performance to an SoC with only 1 of the larger cores. (Should)
    \item Run Linux on the previously designed SoC and connect to it via SSH (Should).
    \item Allow processes to execute on both cores inside of Linux - SMP (Could).
    \item Intelligently select which core the process will run on, depending on factors such as process priority and resource usage (Could).
\end{enumerate}

\clearpage
\section{Risks and Ethical Considerations}
\begin{center}
    \begin{table}[htbp!]
        \begin{tabularx}{\textwidth}{sbbs}
            \hline
            Risk & Risk Description & Mitigation & Risk level \\
            \hline
            FPGA is too small & As the design for the SoC grows larger and more complex, it will require a greater amount of LUTs. This may result in a larger design than the FPGA to be used can implement. & This risk could be mitigated by purchasing a larger FPGA, or by reducing the complexity of the design and objectives. & High \\
            \hline
            Student is unable to work & Either by illness or other matters, the student is unable to continue work on the project for an extended period. & Mitigating this risk is difficult as illness cannot be predicted, but the likelihood of this occurring is very low. & Low \\
            \hline
            Concept is flawed/ impossible/ too difficult & As the project progresses, it appears that the main objectives of the project cannot be completed either at all, or in the time-frame provided. & The objectives and scope of the project can be changed/reduced in order to produce a full piece of work by the deadline. & Low \\
            \hline
        \end{tabularx}
    \end{table}
\end{center}

There are no ethical considerations for this project.
\clearpage

\section{Project Management and Resources}
Weekly meetings will be held with the Project Supervisor to discuss current progress and any issues that may arise during the project. A project work board will be made and kept updated with tasks to complete, completed work and a log of what is currently in progress. This will be shared with the Project Supervisor and assist in tracking the project over it's lifetime, as well as helping in the creation of the Progress Report.

An Artix A7 FPGA prototyping board is the target FPGA to be used during development, and where the final implementation of the SoC will be tested. \texttt{git} will be used as version control for the continuous development of the hardware description, along with \texttt{GitHub} to store a remote copy of all work completed to ensure the risk of losing progress by hardware malfunction is minimal.
\clearpage
\subsection{Timetable}
\begin{center}
    \begin{table}[htbp!]
        \begin{tabularx}{\textwidth}{sb}
            \hline
            Date & Objective \\
            \hline
            13/10/22 & Finalisation and submission of Project Specification \\
            \hline
            13/10/22 - 21/10/22 & Complete research and testing on the capabilities of Rocket Chip generation \\
            \hline
            21/10/22 - 20/11/22 & Create initial design of the heterogeneous CPU and SoC \\
            \hline
            20/11/22 - 10/12/22 & Finalise design of CPU and SoC, begin testing and attempt Linux boot \\
            \hline
            20/11/22 - 28/11/22 & Write and then submit the combined Project Specification and Progress Report \\
            \hline
            10/12/22 - 15/02/23 & Complete testing and attempted Linux boot, potential research into applications of the designed SoC \\
            \hline
            15/02/23 - 6/03/23 & Write project presentation \\
            \hline
            6/03/23 & Present the project outcomes and demonstrate work \\
            \hline
            6/03/23 - 15/04/23 & Completion of first draft of Final Report \\
            \hline
            15/04/23 - 2/05/23 & Finalisation and submission of the Final Report \\
            \hline
        \end{tabularx}
    \end{table}
\end{center}

% \chapter{Introduction}
\label{ch:introduction}

The aim of this project is to design a heterogeneous system on a chip (SoC) using the RISC-V open standard ISA\cite{riscv-1}\cite{riscv-2}. Once designed, the SoC should be implementable on an FPGA and be able to execute bare-metal C code and assembly with SMP, allowing all cores to execute code at the same time.

\section{Motivation}
Designing a processor is a complex challenge. They need to be power efficient to reduce the cost of running the system, as well as providing adequate processing power when required. Often one is sacrificed to improve the other: for instance, energy efficiency is often reduced due in designs of large processors containing multiple cores. Due to the increased size and complexity, more power has to be supplied to achieve the same performance as a smaller processor.

A heterogeneous CPU attempts to solve this issue by combining dissimilar core designs; often a powerful core, or p-core, and an efficient core, or e-core. When only light processing is required, the p-core can be effectively shutdown and the e-core will do all processing, resulting in less power used. For heavy processing, the p-core is then used to increase peak performance. Depending on the exact implementation, the p-core can be used either individually or in tandem with the e-core, but both result in greater performance than just the e-core.

This has been used extensively in many mobile devices, and increasingly in larger devices like laptops. The Apple M2 processor\cite{applem2} is one such example, implementing 4 high-performance cores and 4 energy-efficient cores. These are arranged in a hybrid configuration similar to ARM DynamIQ and big.LITTLE\cite{biglittle}, allowing cores to dynamically be assigned work that best fits them. In Apple's design, both cores are capable of executing the same code, with the performance core benefitting from much increased cache and other proprietary changes to increase it's speed over the efficiency core. 

\section{Objectives}
%One sentence summary of your project. Followed by a short list of concrete objectives:
The overall aim is design a heterogeneous SoC containing two types of RISC-V core that can execute bare-metal C code and assembly with SMP. Objectives have been labelled according to the MoSCoW method\cite{case-method-fasttrack} to indicate their importance and expected completion.
\begin{figure}[h!]
    \centering
    \begin{enumerate}
        \item Design a heterogeneous RISC-V SoC containing 2 dissimilar cores, each capable of executing at least the RV64 instruction set (Must).
        \item Execute assembly instructions on both cores when implemented in an SoC on an FPGA (Must).
        \item Execute bare-metal code on both cores when implemented in an SoC on an FPGA (Should).
        \item Measure the performance of the SoC for comparison against single-core designs. (Should)
        \item Run embedded Linux on the SoC and connect to it via serial or SSH (Could).
        \item Allow processes to execute on both cores inside of embedded Linux (Could).
        \item Intelligently select which core the process will run on, depending on factors such as process priority and resource usage (Won't).
    \end{enumerate}
    \caption{Objectives for the project}
    \label{fig:objectives}
\end{figure}
% \chapter{Background and Research}
\label{ch:background}

\section{RISC-V}
\label{sec:riscv-background}
\subsection{ISAs and RISC vs CISC}
An ISA is an instruction set architecture, and is the formal definition for what an abstract model of a CPU. The ISA defines the instructions, registers, communication standards and other features of a CPU in order to allow a person with the ISA to design a physical CPU that would implement the abstract model. CPUs that implement the same ISA (or supersets of an ISA) are able to execute the same code, meaning programs written for a CPU can be run on a different CPU if they share the same ISA. This is how a single compiled program is able to be run on so many different systems, and is incredibly important for modern devices, where there is a huge range of processors available. 

Many modern ISAs are based on older ISAs that have been extended to allow for new instructions, registers, communication protocols, etc, while maintaining support for previous ISAs. This is typically called CISC (Complex Instruction Set Computers), where there are a very large amount of instructions the processor can execute. The backwards compatibility can be very useful, and having explicit instructions for many tasks can mean the tasks are completed more efficiently than if multiple smaller instructions were executed. Having large instructions can also reduce program sizes, due to the increased operations per instruction, which is very beneficial, especially in small systems.

There are several issues with CISC however, such as increased complexity in the CPU as the amount of instructions increases, leading to rising costs of design and manufacture, and can force the physical size of the CPU to be larger to accommodate the extra logic space for larger instructions. More complex designs will also make future development harder if backwards compatibility is to be maintained. 

The alternative to CISC is RISC (Reduced Instruction Set Computer). RISC aims to reduce the number of instructions the CPU requires in order to keep the design as simple as possible. This allows the design to be done much more easily, as well as allowing future extensions to be easier. Reduced complexity can also allow for greater optimisations in what instructions are in the ISA, with the aim of achieving greater performance than CISC by executing more instructions at much greater speed. There are drawbacks however, with increased program size and possible reduction in speed of execution for multiple operations that could've been completed in a single instruction using CISC.

\subsection{The RISC-V ISA}
RISC-V is a new, modern, open-standard ISA that follows RISC concepts. The ISA does not specify a single instruction set, but multiple by having three different width instruction sets with different base instructions. 32-bit RISC-V (RV32I) is focussed on embedded and personal systems, 64-bit RISC-V (RV64I) for personal and server systems and 128-bit RISC-V (RV128I) for server and high-performance compute systems. RV32 and RV64 have embedded versions, RV32E and RV64E, that have reduced registers and instructions.

\subsubsection{Extensions}
At base, these implement only integer addition/subtraction. The ISA contains extensions that add more functionality to the CPUs, such as support for multiplication/division, floating point operations IEEE-754, compressed instructions, atomics, etc. An instruction set implementing these extensions is referred to by adding the initial of the extension to the name, so a 64-bit CPU with integer addition/subtraction, multiplication/division and compressed instructions would be `RV64IMC`. Custom extensions can also be written, allowing for a great deal of customisation in RISC-V CPU designs.

There are 30 official extensions to the base instruction sets. The most popular are listed below.
\begin{description}
    \item[M] Integer multiplication and division
    \item[A] Atomic instructions
    \item[F] Single-precision floating point
    \item[D] Double-precision floating point
    \item[Zicsr] Control and Status Register (CSR), useful in multi-CPU systems
    \item[Zifencei] Instruction-Fetch Fence
    \item[G] Short for IMAFDZicsrZifencei as a 'general-purpose' CPU, e.g. RV64G
    \item[C] Compressed instructions, 16-bit instead of 32/64
    \item[Others] Quad-precision floating point, bit manipulation, vector, misaligned atomics, etc
\end{description}

\section{Processor Design}
Processor performance can be measured in multiple ways, with the main metrics being speed at executing tasks, cost of manufacturing, physical size and energy consumption. These are implicitly linked - a processor with larger physical size will likely draw a larger amount of energy, contain more logic and so be faster at executing tasks and cost more to manufacture. The aim of processor design is to maximise the speed of task execution, while minimising the rest. The balance between these is what gives rise to different processor designs, as a mobile device is much more constrained in power usage than a large server and will require a different processor.

Maximising task execution while minimising energy consumption can be balanced in a multitude of ways. 

\section{Related work}
%TODO: To be updated - mostly copied from specification, want this to be much more critical/analytical
%Discuss related work.
\subsection{A RISC-V Heterogeneous SoC for Embedded Devices\cite{valenterisc}}
This project is ongoing, and presents work designing a RV64 (RISC-V 64-bit) host core that offloads tasks to a PMCA (Programmable Many Core Accelerator) made from RV32 (RISC-V 32-bit) cores, which implement extensions for machine learning and discrete signal processing. The suggested use-case for the SoC is in IoT applications and programmable embedded devices. The host core is Linux compatible, and offers a full OS that acts as a platform for programs that run on the PMCA. The use of a large RV64 core to allow a full Linux OS to run on the SoC provides a huge amount of flexibility to the programmer, as the OS implements features like CLI, memory virtualisation, networking and more that allow programs to be written much more generally than embedded software running without an OS. However, this usage of a full Linux OS could be considered excessive for the use-case. An embedded Linux OS would have a lower overhead due to the reduced services it offers, which is very beneficial in an embedded environment where efficiency is highly important. Unfortunately, no data is provided about the processing power or energy usage of the design.

\subsection{Muntjac multicore RV64 processor\cite{UCAM-CL-TR-972}}
Muntjac is an SoC generator comprising of multiple components, that can be used to produce a Linux capable SoC. There is only one type of core in the system, RV64, but the SoC can be multicore. The project report is dedicated to the design of the core and cache as opposed to usage in any devices, as the purpose is to provide an easily understood and extensible platform for specialised designs. This is excellently done - the project is very well presented and uses multiple open-standards like TileLink\cite{tilelink} to increase the ease of working with it. It would be possible for this project to be extended for a heterogeneous SoC design, but it appears that this has not yet been done in any public works that extend the project.
% \section{Progress}
\label{sec:progress}

Summarise the progress you have made so far. You can cross-reference other sections (\Cref{sec:background}).
% \section{Project management}

Include a timetable (in 2 week chunks) for the remainder of the academic year, up until the submission deadline.

\bibliographystyle{ieeetr}
\bibliography{../common/bibliography}

\end{document}